\documentclass[10pt,a4paper]{article}
\usepackage[margin=2cm]{geometry}
\usepackage[T1]{fontenc}
\usepackage{lmodern}
\usepackage{microtype}
\usepackage{amsmath,amssymb}
\makeatletter\renewcommand\verbatim@font{\footnotesize\ttfamily}\makeatother
\sloppy
\usepackage{booktabs}
\usepackage{tabularx}
\usepackage{longtable}
\usepackage{enumitem}
\PassOptionsToPackage{hyphens}{url}
\usepackage[hidelinks]{hyperref}
\setlength{\emergencystretch}{3em}
\setlist{nosep}
\newlist{steps}{itemize}{1}
\setlist[steps]{label=$\square$, leftmargin=1.5em, itemsep=2pt}
\newcolumntype{Y}{>{\raggedright\arraybackslash}X}
\newcommand{\f}[1]{\texttt{#1}}

\title{Implementation Plan\\[2pt]\large Part 1 as a study: the PhonePe Pulse state panel, 2018--2026}
\date{4 November 2026}
\begin{document}
\maketitle

\noindent\textbf{For agentic workers:} REQUIRED SUB-SKILL: use \f{superpowers:subagent-driven-development}
(recommended) or \f{superpowers:executing-plans} to implement this plan task by task. Each step is a
checkbox ($\square$); mark it done ($\blacksquare$) when complete.

\paragraph{Goal.} Replace the two-page Part 1 progress report with a 20--30 page study,
\f{thesis/study/part1.tex}, in paper prose with STE words, where every number comes from code.

\paragraph{Architecture.} Python adds one PhonePe input (transaction count by type) and a data-check file.
One R script reads the Part 1 interim files and writes the study figures (PDF), tables (\f{.tex}) and a file
of computed numbers (\f{numbers.tex}). The document prints each number with \verb|\val{key}|; an unknown
key stops the compile.

\paragraph{Tech stack.} Python 3.12 (pandas, pytest) in \f{.venv}; R 4.3.3 with \f{renv} (data.table,
ggplot2, testthat); pdflatex, BibTeX (\f{plainnat}), \f{latexmk}; GNU Make; Git. No new packages.

\paragraph{Spec.} \f{docs/superpowers/specs/2026-11-04-part1-study-design.tex}. Read it before each task.
Background: \f{docs/superpowers/specs/2026-11-01-upi-financial-inclusion-design.tex},
\f{thesis/progress/part1.tex}.

\section*{Global constraints}
\begin{itemize}
\item Branch \f{btp1}. Do not push. Do not move or delete tags. Do not rewrite history. Do not backdate.
\item Each commit message ends with the line \f{Co-Authored-By: Claude Opus 5.5 <noreply@anthropic.com>}
      (after one blank line).
\item Do not change these files: \f{data/interim/phonepe.csv}, \f{data/interim/population.csv},
      \f{data/interim/cpi.csv}, \f{data/interim/npci.csv}, \f{data/final/*}, \f{output/tables/*},
      \f{src/py/btpdata/states.py}, \f{state\_crosswalk.csv}, \f{src/R/01\_descriptives.R},
      \f{thesis/btp1/*}, \f{thesis/btp1-slides/*}, \f{thesis/progress/*}.
\item PhonePe commit \f{257bf38e3095020f9f516c4c72c10bdc348781b1}. Periods \f{2018Q1}--\f{2026Q2} (34).
      35 \f{state\_id} values; main sample 30 (\f{btpdata.states.MAIN\_SAMPLE}).
\item Interim schema: \f{state\_id,period,variable,value,source}, in this order.
\item Python modules run from the repo root: \f{PYTHONPATH=src/py .venv/bin/python -m btpdata.<module>}.
      R scripts run from the repo root: \f{Rscript src/R/<script>.R}. R tests run with
      \f{Rscript -e 'testthat::test\_dir("tests/R")'} (working directory \f{tests/R}, so tests use
      \f{../../} paths).
\item Scope: no CPI, no real values, no NPCI numbers, no maps, no $\beta$-regression (spec Section 2).
\item Statistics across states use the 30 main-sample ids. Figures and tables show all 35 ids and mark the
      five other union territories. Reference quarters: 2018Q4 and 2026Q2.
\item Units: transactions in billion ($10^9$) or per person per quarter; amounts in Rs.\ or Rs.\ lakh crore
      ($10^{12}$ rupees); users in million or per 1000 persons; population in million; shares in percent.
\item Document text: ASD-STE100 vocabulary and sentence rules (common words, short sentences, active voice),
      in paragraphs. No filler, idioms, hype words or value judgments. No reference to ``the student'', to
      AI, or to plan task numbers. The author is Devansh Gupta.
\item Every number in the document prose comes from \verb|\val{key}| or from an \verb|\input| table. File
      paths, the commit hash, dates of access and section numbers are not numbers in this sense.
\item New sources: one row in \f{references/source-log.csv} and one entry in \f{references/refs.bib}. PDFs
      only in \f{references/pdfs/} (ignored by Git).
\end{itemize}

\section*{Review focus}
\begin{enumerate}
\item A type entry is absent from a state-quarter file (real case: Lakshadweep 2018Q1 \f{Utility}). Expect a
      zero only when the other types sum to the hover count; else a stop that names the file.
      Tested in Task~2 (\f{test\_zero\_fill}, \f{test\_absent\_type\_bad\_sum\_raises}).
\item The raw clone already exists at the pinned commit with the old sparse set. Expect \f{fetch} to add the new
      paths, not to return early. Tested in Task~1 (\f{test\_fetch\_extends\_sparse\_set}).
\item Merged ids in the type data: Ladakh and Jammu and Kashmir are separate folders. Expect one
      \f{jammu\_kashmir} row per type that is the sum. Tested in Task~2 (\f{test\_values}).
\item A key used in the document has no value, or a value is \f{NA}/\f{Inf}. Expect the build to stop. Tested
      in Task~4 (\f{test\_non\_finite}) and Task~6 (compile with an unknown key).
\item Ties in ranks (equal values across states). Expect average ranks in the Spearman coefficient, as in the
      Part 2 validation. Tested in Task~3 (\f{test\_spearman\_ties}).
\end{enumerate}

\section*{File map}
\par\noindent\begin{tabularx}{\textwidth}{@{}>{\raggedright\arraybackslash}p{0.38\textwidth} Y@{}}
\toprule
File & Responsibility \\
\midrule
\f{src/py/btpdata/fetch\_phonepe.py} (modify) & sparse paths; re-apply a changed sparse set \\
\f{src/py/btpdata/} \f{parse\_phonepe\_types.py} & count by type; checks $\to$ \f{data/interim/phonepe\_types.csv},
  \f{phonepe\_checks.csv} \\
\f{src/R/lib/state\_labels.csv} & \f{state\_id,name,code,in\_main} for the 35 ids \\
\f{src/R/lib/part1\_fns.R} & loading, measures, national series, concentration, decomposition, ranks \\
\f{src/R/lib/latex\_out.R} & number formatting, \f{numbers.tex} writer, table writer \\
\f{src/R/lib/study\_theme.R} & ggplot2 theme, period axis, PDF save \\
\f{src/R/part1\_study.R} & all study outputs under \f{output/study/part1/} \\
\f{thesis/study/part1.tex}, \f{part1.pdf} & the study \\
\f{tests/py/test\_fetch\_phonepe.py}, \f{test\_phonepe\_types.py}, \f{test\_state\_labels.py} & Python tests \\
\f{tests/py/fixtures/phonepe\_types/} & fixture repository for the type parser \\
\f{tests/R/test\_part1.R}, \f{test\_latex\_out.R} & R tests \\
\f{Makefile} (modify) & \f{data}, \f{analysis}, \f{report} targets \\
\bottomrule
\end{tabularx}

% =====================================================================
\section{Task 1: Fetch the new PhonePe paths}
\textbf{Files.} Modify \f{src/py/btpdata/fetch\_phonepe.py}. Create \f{tests/py/test\_fetch\_phonepe.py}.
Changes \f{data/raw/phonepe/manifest.csv} (committed) and the ignored clone.

\textbf{Interfaces (produces).}
\f{SPARSE\_PATHS: list[str]} = the two current paths plus \f{data/aggregated/transaction/country/india} and
\f{data/aggregated/user/country/india}.
\f{fetch(dest: Path, commit: str = PHONEPE\_COMMIT, paths: list[str] | None = None, url: str = PHONEPE\_URL) -> Path}
(\f{paths=None} means \f{SPARSE\_PATHS}).
Behaviour: fresh clone as now, then \f{sparse-checkout set *paths}, checkout, manifest row. If the clone is at
\f{commit} and \f{set(git sparse-checkout list) == set(paths)}: return, no manifest row. If the clone is at
\f{commit} and the sets differ: \f{sparse-checkout set *paths}, then a manifest row. Manifest \f{notes}:
\f{"git commit <commit>; sparse paths: <p1>;<p2>;..."}.

\begin{steps}
\item Write the failing tests. A fixture function makes a local Git repository in \f{tmp\_path/"src"} with one
      file under each of the four paths (for example \f{<path>/2018/1.json} with \f{\{\}}) and one commit.
\begin{verbatim}
def test_fetch_extends_sparse_set(tmp_path):
    src, sha = make_repo(tmp_path)
    old = SPARSE_PATHS[:2]
    repo = fetch(tmp_path / "raw", commit=sha, paths=old, url=str(src))
    assert (repo / old[0] / "2018/1.json").exists()
    assert not (repo / "data/aggregated/transaction/country/india/2018/1.json").exists()
    fetch(tmp_path / "raw", commit=sha, paths=SPARSE_PATHS, url=str(src))
    assert (repo / "data/aggregated/transaction/country/india/2018/1.json").exists()
    rows = list(csv.DictReader(open(tmp_path / "raw/manifest.csv")))
    assert len(rows) == 2 and "data/aggregated/user/country/india" in rows[1]["notes"]
def test_fetch_same_set_no_new_row(tmp_path):
    src, sha = make_repo(tmp_path)
    fetch(tmp_path / "raw", commit=sha, paths=SPARSE_PATHS, url=str(src))
    fetch(tmp_path / "raw", commit=sha, paths=SPARSE_PATHS, url=str(src))
    assert len(list(csv.DictReader(open(tmp_path / "raw/manifest.csv")))) == 1
def test_sparse_paths(): assert len(SPARSE_PATHS) == 4
\end{verbatim}
\item Run \f{.venv/bin/pytest tests/py/test\_fetch\_phonepe.py -v}. Expected: FAIL (\f{fetch} has no
      \f{paths}/\f{url} argument).
\item Implement. Use \f{\_git(repo, "sparse-checkout", "list").splitlines()} for the current set.
\item Run the new tests and \f{.venv/bin/pytest -q}. Expected: all PASS.
\item Run on real data: \f{PYTHONPATH=src/py .venv/bin/python -m btpdata.fetch\_phonepe}. Expected: exit 0;
      \f{ls data/raw/phonepe/repo/data/aggregated/transaction/country/india/state | wc -l} prints 36;
      \f{data/raw/phonepe/manifest.csv} has one new row.
\item Check: \f{PYTHONPATH=src/py .venv/bin/python -m btpdata.parse\_phonepe \&\& git diff --exit-code data/interim/phonepe.csv}.
      Expected: exit 0.
\item Commit: \f{git add src/py/btpdata/fetch\_phonepe.py tests/py/test\_fetch\_phonepe.py data/raw/phonepe/manifest.csv}
      and \f{git commit -m "feat: fetch PhonePe aggregated transaction and user paths"}.
\end{steps}

% =====================================================================
\section{Task 2: Transaction count by type, and data checks}
\textbf{Files.} Create \f{src/py/btpdata/parse\_phonepe\_types.py}, \f{tests/py/test\_phonepe\_types.py},
\f{tests/py/fixtures/phonepe\_types/data/...}. Modify \f{Makefile} (\f{data} target: add
\f{\$(MOD).parse\_phonepe\_types} after \f{\$(MOD).parse\_phonepe}). Produces \f{data/interim/phonepe\_types.csv},
\f{data/interim/phonepe\_checks.csv}.

\textbf{Interfaces.} Consumes \f{parse\_phonepe.\_load}, \f{parse\_phonepe.HOVER}, \f{parse\_phonepe.SOURCE},
\f{states.to\_state\_id}, \f{states.sum\_merged}, \f{periods.quarters}. Produces:
\begin{itemize}
\item \f{TYPES = \{"P2P": "phonepe\_count\_p2p", "Retail": "phonepe\_count\_retail", "Utility": "phonepe\_count\_utility"\}}.
\item \f{folder\_to\_raw(folder: str) -> str}: \f{folder.replace("-", " ")}. The crosswalk does not change.
\item \f{parse\_type\_file(path: Path) -> dict[str, int]}: type name $\to$ count of the \f{TOTAL} instrument;
      a name outside \f{TYPES} raises \f{ValueError} with the name and the path.
\item \f{build\_types(repo: Path, out: Path, checks\_out: Path, periods: list[str]) -> tuple[pd.DataFrame, pd.DataFrame]}.
\item \f{main()}: calls \f{build\_types} with the repository \f{data/raw/phonepe/repo}, the outputs
      \f{data/interim/phonepe\_types.csv} and \f{data/interim/phonepe\_checks.csv}, and \f{quarters()}.
\end{itemize}

\textbf{Rules (spec Section 4.2).} For each period: the state folders are the subfolders of
\f{data/aggregated/transaction/country/india/state/}. Each folder maps (via \f{folder\_to\_raw}) to one hover name of
\f{map/transaction/hover} for that period; a hover name with no folder, or a folder with no hover name, raises
\f{ValueError}. A missing type file raises \f{FileNotFoundError} (the path is in the message). An absent type is
set to 0 only if the present types sum to the hover count; else \f{ValueError} naming the file. If the full sum
differs from the hover count: \f{ValueError} naming the file. The national file
\f{aggregated/transaction/country/india/<y>/<q>.json} (sum of its type counts) must equal the sum of hover counts,
and \f{aggregated/user/country/india/<y>/<q>.json} (\f{registeredCount}) must equal the sum of
\f{map/user/hover} counts; else \f{ValueError} naming the file. Then \f{sum\_merged} on the raw names, keys
\f{["period","variable"]}; \f{source = SOURCE}.

\textbf{Checks file.} Columns \f{check,n\_items,n\_flagged,max\_abs\_diff,note}, rows in this order:
\par\noindent\begin{tabularx}{\textwidth}{@{}>{\raggedright\arraybackslash}p{0.38\textwidth} Y@{}}
\f{type\_sum\_equals\_state\_total} & items: raw state-quarters; flagged 0 \\
\f{absent\_type\_set\_to\_zero} & items: raw state-quarter-type slots ($3\times$ state-quarters); flagged: slots
  set to 0; note: \f{"<folder> <period> <type>"} joined by \f{"; "} \\
\f{national\_count\_equals\_state\_sum} & items: quarters; flagged 0 \\
\f{national\_users\_equals\_state\_sum} & items: quarters; flagged 0 \\
\end{tabularx}

\begin{steps}
\item Write the fixture repository \f{tests/py/fixtures/phonepe\_types/} for 2018Q1 (JSON layout as in the pinned
      commit; see \f{parse\_phonepe.py} and the spec). Hover transactions: andhra pradesh (count 10862829,
      amount 14001071206.72507), jammu \& kashmir (100, 1000.0), ladakh (10, 50.0). Hover users: 500, 40, 5.
      Type files in folders \f{andhra-pradesh} (P2P 4000000, Retail 6000000, Utility 862829),
      \f{jammu-\&-kashmir} (40, 50, 10), \f{ladakh} (6, 4, no Utility entry). National transaction file: P2P
      4000046, Retail 6000054, Utility 862839. National user file: 545.
\item Write the failing tests:
\begin{verbatim}
def test_folder_to_raw(): assert to_state_id(folder_to_raw("jammu-&-kashmir")) == "jammu_kashmir"
def test_values(tmp_path):
    df, chk = build_types(FIX, tmp_path/"t.csv", tmp_path/"c.csv", ["2018Q1"])
    v = df.set_index(["state_id", "variable"]).value
    assert v[("andhra_pradesh", "phonepe_count_utility")] == 862829
    assert v[("jammu_kashmir", "phonepe_count_p2p")] == 46
    assert v[("jammu_kashmir", "phonepe_count_retail")] == 54
    assert v[("jammu_kashmir", "phonepe_count_utility")] == 10
    assert list(df.columns) == ["state_id", "period", "variable", "value", "source"]
def test_zero_fill(tmp_path):
    _, chk = build_types(FIX, tmp_path/"t.csv", tmp_path/"c.csv", ["2018Q1"])
    z = chk.set_index("check").loc["absent_type_set_to_zero"]
    assert z.n_items == 9 and z.n_flagged == 1 and z.note == "ladakh 2018Q1 Utility"
    assert chk.check.tolist() == ["type_sum_equals_state_total", "absent_type_set_to_zero",
                                  "national_count_equals_state_sum", "national_users_equals_state_sum"]
def test_sum_mismatch_raises(tmp_path):     # copy FIX; AP P2P -> 4000001
    with pytest.raises(ValueError, match="andhra-pradesh/2018/1.json"): ...
def test_absent_type_bad_sum_raises(tmp_path):  # copy FIX; ladakh Retail -> 3
    with pytest.raises(ValueError, match="ladakh/2018/1.json"): ...
def test_national_mismatch_raises(tmp_path):    # copy FIX; national P2P -> 4000047
    with pytest.raises(ValueError, match="india/2018/1.json"): ...
def test_unknown_type_raises(tmp_path):         # copy FIX; AP type name "Other"
    with pytest.raises(ValueError, match="Other"): ...
@pytest.mark.skipif(not REAL.exists(), reason="raw clone absent")
def test_real_types_sum_to_count():
    t = pd.read_csv("data/interim/phonepe_types.csv"); p = pd.read_csv("data/interim/phonepe.csv")
    assert len(t) == 3570 and t.state_id.nunique() == 35
    s = t.groupby(["state_id", "period"]).value.sum()
    c = p[p.variable == "phonepe_count"].set_index(["state_id", "period"]).value
    assert (s - c.loc[s.index]).abs().max() == 0
\end{verbatim}
\item Run \f{.venv/bin/pytest tests/py/test\_phonepe\_types.py -v}. Expected: FAIL (ImportError).
\item Implement \f{parse\_phonepe\_types.py}.
\item Run the tests. Expected: PASS except \f{test\_real\_types\_sum\_to\_count} (file not yet written).
\item Run on real data: \f{PYTHONPATH=src/py .venv/bin/python -m btpdata.parse\_phonepe\_types}. Expected: exit 0;
      \f{data/interim/phonepe\_checks.csv} gives \f{type\_sum\_equals\_state\_total} 1224/0,
      \f{absent\_type\_set\_to\_zero} 3672/1 with note \f{lakshadweep 2018Q1 Utility}, and 34/0 for the two
      national checks. If a check fails, stop and report the file and the values. Do not change a rule.
\item Run \f{.venv/bin/pytest -q}. Expected: all PASS.
\item Run \f{git diff --exit-code data/interim/phonepe.csv data/interim/population.csv data/final}. Expected: exit 0.
\item Commit the module, tests, fixtures, \f{Makefile}, the two new interim files:
      \f{git commit -m "feat: PhonePe transaction count by type, with data checks"}.
\end{steps}

% =====================================================================
\section{Task 3: State labels and the R measures library}
\textbf{Files.} Create \f{src/R/lib/state\_labels.csv}, \f{tests/py/test\_state\_labels.py},
\f{src/R/lib/part1\_fns.R}, \f{tests/R/test\_part1.R}. Modify \f{references/source-log.csv},
\f{references/refs.bib}.

\textbf{state\_labels.csv.} Columns \f{state\_id,name,code,in\_main}; 35 rows sorted by \f{state\_id};
\f{in\_main} is \f{TRUE}/\f{FALSE}. Names: the official name; \f{jammu\_kashmir} is ``Jammu and Kashmir (with
Ladakh)''; \f{dnh\_dd} is ``Dadra and Nagar Haveli and Daman and Diu''; \f{andaman\_nicobar} is ``Andaman and
Nicobar Islands''. Codes: the ISO 3166-2:IN subdivision code without \f{IN-}. Expected codes (check each against
the source and use the source value): AN AP AR AS BR CH CG DH DL GA GJ HR HP JK JH KA KL LD MP MH MN ML MZ NL OD
PB PY RJ SK TN TG TR UP UK WB. Source: the Wikipedia page ``ISO 3166-2:IN''
(\url{https://en.wikipedia.org/wiki/ISO_3166-2:IN}); log it as \f{D11} in \f{source-log.csv} (type
\f{dataset}, used for ``two-letter state codes for figure labels'') and as \f{@misc\{iso3166\_2\_in, ...\}} in
\f{refs.bib} with \f{srcid = \{D11\}} and the access date.

\textbf{Interfaces (produces, \f{part1\_fns.R}).} All data are \f{data.table}. \f{periods} are \f{"YYYYQn"}.
\begin{itemize}
\item \f{load\_part1(phonepe, types, population, labels) -> data.table}: one row per \f{state\_id}, \f{period}
      (1190 rows on real data), columns \f{state\_id, period, name, code, in\_main, count, amount, users,
      population, count\_p2p, count\_retail, count\_utility}. Stops if a state-period is missing in any input.
\item \f{add\_measures(dt) -> dt} adds \f{count\_pc} ($C/P$), \f{amount\_pc} ($A/P$), \f{users\_per\_1000}
      ($1000U/P$), \f{intensity} ($C/U$), \f{atv} ($A/C$), \f{share\_p2p}, \f{share\_retail},
      \f{share\_utility} ($C^k/C$, fractions).
\item \f{national(dt) -> data.table}: by period, sums of \f{count, amount, users, population, count\_*}, then the
      same measures as \f{add\_measures}, plus \f{qoq\_count}, \f{qoq\_amount} (percent change on the previous
      quarter; \f{NA} for 2018Q1).
\item \f{annual(nat) -> data.table}: one row per calendar year; \f{year} (character; \f{"2026 (Q1--Q2)"} for the
      last), \f{count} and \f{amount} (sums), \f{users} (last quarter of the year), \f{count\_pc} (count / mean
      population), \f{atv}, \f{growth\_count} (percent on the previous full year; \f{NA} for 2018 and 2026).
\item \f{years\_between(start, end) -> numeric}: (quarter index of end $-$ quarter index of start) / 4.
\item \f{annual\_rate(x0, x1, years) -> numeric}: $100\,((x_1/x_0)^{1/\text{years}} - 1)$.
\item \f{spread(dt, var, ids) -> data.table}: by period, \f{med}, \f{q1}, \f{q3} (R type-7 quantiles) of \f{var} over
      \f{ids}.
\item \f{top\_k(dt, per, k = 5) -> data.table}: the \f{k} ids with the largest \f{count} in \f{per} (all 35), with
      columns \f{variable} (\f{count, amount, users, population}) and \f{share} (percent of the 35-id total
      held by these \f{k} ids).
\item \f{n\_for\_share(x, s = 0.5) -> integer}: the smallest number of units, largest first, whose sum is at least
      \f{s} of the total.
\item \f{spearman(x, y) -> numeric}: Pearson correlation of \f{rank(x)} and \f{rank(y)} (average ranks).
\item \f{decompose(dt, per, ids) -> list}: \f{n}, \f{var\_y}, \f{var\_r}, \f{var\_i}, \f{cov\_ri}, \f{reach},
      \f{intensity} with $y=\ln(C/P)$, $r=\ln(U/P)$, $i=\ln(C/U)$ over \f{ids} in \f{per}; sample variances
      ($n-1$); $\text{reach} = \operatorname{cov}(r,y)/\operatorname{var}(y)$,
      $\text{intensity} = \operatorname{cov}(i,y)/\operatorname{var}(y)$.
\item \f{decompose\_series(dt, ids) -> data.table}: \f{decompose} for each period, one row per period.
\end{itemize}

\begin{steps}
\item Write \f{state\_labels.csv} and log the source in both files.
\item Write \f{tests/py/test\_state\_labels.py}: the \f{state\_id} set equals \f{ALL\_IDS}; the \f{in\_main}
      \f{TRUE} set equals \f{MAIN\_SAMPLE}; codes are two upper-case letters and unique. Run it. Expected: PASS.
\item Write the failing R tests in \f{tests/R/test\_part1.R} (source \f{../../src/R/lib/part1\_fns.R}):
\begin{verbatim}
test_that("decompose terms sum to 1", {   # 4 states, C/U/P chosen by hand
  d <- data.table(state_id = letters[1:4], period = "2020Q1", count = c(10, 40, 90, 160),
                  users = c(5, 10, 15, 20), population = c(100, 100, 100, 100))
  r <- decompose(d, "2020Q1", letters[1:4])
  expect_equal(r$reach + r$intensity, 1); expect_equal(r$n, 4L)
  expect_equal(r$var_y, var(log(d$count / 100)))
})
test_that("top_k", {   # counts 50, 30, 10, 10; population 1, 1, 1, 7
  ...; expect_equal(t[variable == "count"]$share, 80); expect_equal(t[variable == "population"]$share, 20)
})
test_that("n_for_share", { expect_equal(n_for_share(c(50, 30, 10, 10)), 1L)
                           expect_equal(n_for_share(c(40, 30, 20, 10)), 2L) })
test_that("spearman_ties", { expect_equal(spearman(c(1, 2, 2, 3), c(1, 2, 3, 4)),
                                          cor(c(1, 2.5, 2.5, 4), 1:4)) })
test_that("annual_rate", { expect_equal(annual_rate(100, 121, 2), 10)
                           expect_equal(years_between("2018Q4", "2026Q2"), 7.5) })
test_that("national qoq", { # 2 states; count 10 and 10 in 2020Q1, 15 and 15 in 2020Q2
  expect_equal(n$qoq_count, c(NA, 50)) })
test_that("annual", { 8 quarters of 2018-2019 with count 1 each; expect_equal(a$count, c(4, 4))
                      expect_equal(a$growth_count[2], 0) })
test_that("load_part1 stops on a missing state-period", { ... expect_error(..., "missing") })
\end{verbatim}
\item Run \f{Rscript -e 'testthat::test\_dir("tests/R")'}. Expected: the new tests FAIL (functions not found).
\item Implement \f{part1\_fns.R}.
\item Run the R tests and \f{.venv/bin/pytest -q}. Expected: all PASS.
\item Commit: \f{git commit -m "feat: state labels and Part 1 study measures (R)"}.
\end{steps}

% =====================================================================
\section{Task 4: LaTeX output writers (R)}
\textbf{Files.} Create \f{src/R/lib/latex\_out.R}, \f{tests/R/test\_latex\_out.R}.

\textbf{Interfaces (produces).}
\begin{itemize}
\item \f{fmt(x, digits) -> character}: fixed decimals; thousands separator \verb|{,}| in the integer part; minus as
      \verb|$-$|. Example: \f{fmt(38664.285, 1)} is \verb|38{,}664.3|; \f{fmt(-10.56, 1)} is \verb|$-$10.6|.
\item \f{latex\_escape(x) -> character}: escapes \verb|& % $ # _ { }|.
\item \f{nv(key, value, digits = NULL, source) -> data.table}: one entry; numeric values go through \f{fmt};
      character values through \f{latex\_escape}. Stops if \f{key} does not match
      \verb|^[a-z][a-z0-9_]*(\.[A-Za-z0-9_]+)*$| or if a numeric value is not finite.
\item \f{write\_numbers(entries, path, generator)}: stops on a duplicate key; writes a header comment
      (\f{generated by <generator>; do not edit}) and one line per entry:
      \verb|\defval{<key>}{<value>}  % <source>|, in the order given.
\item \f{write\_tabular(cells, path, align, header, env = "tabular", caption = NULL, label = NULL)}: \f{cells} is a
      character table (already formatted); writes \verb|\begin{<env>}{<align>}|, \verb|\toprule|, \f{header}
      (raw LaTeX, may span lines), \verb|\midrule|, rows joined by \verb|&| and ended by \verb|\\|,
      \verb|\bottomrule|, \verb|\end{<env>}|. For \f{env = "longtable"}: \f{caption} and \f{label} at the top,
      then \verb|\endfirsthead| with the header repeated and \verb|\endhead|. Stops if the column count of
      \f{cells} differs from the count of column letters in \f{align}.
\end{itemize}

\begin{steps}
\item Write the failing tests (\f{test\_latex\_out.R}): \f{fmt} examples above and \f{fmt(0.0000015, 7)} is \f{0.0000015};
      \f{latex\_escape("A \& B\_1")} is \verb|A \& B\_1|; \f{test\_bad\_key} (\f{"Nat.count"} and
      \f{"nat count"} stop); \f{test\_non\_finite} (\f{nv("a.b", NA\_real\_, 1, "x")} and \f{Inf} stop);
      \f{test\_duplicate} (two entries with one key stop); \f{test\_numbers\_file} (written file has the
      header line and \verb|\defval{a.b}{1.5}  % src|); \f{test\_tabular} (two rows, \f{align = "lr"}: file
      text equals the expected string; three cells with \f{"lr"} stops).
\item Run \f{Rscript -e 'testthat::test\_dir("tests/R")'}. Expected: new tests FAIL.
\item Implement \f{latex\_out.R} with base R only.
\item Run the R tests. Expected: all PASS.
\item Commit: \f{git commit -m "feat: LaTeX number and table writers (R)"}.
\end{steps}

% =====================================================================
\section{Task 5: Study script: figures, tables, numbers}
\textbf{Files.} Create \f{src/R/lib/study\_theme.R}, \f{src/R/part1\_study.R}. Modify \f{Makefile}
(\f{analysis}: add \f{Rscript src/R/part1\_study.R}). Produces \f{output/study/part1/\{figures,tables\}/*},
\f{output/study/part1/numbers.tex}.

\textbf{Interfaces.} Consumes Tasks 3 and 4. \f{study\_theme.R} produces \f{theme\_study()} (the theme of
\f{01\_descriptives.R}: base size 9, sans, palette \f{\#0b0b0b}, \f{\#52514e}, grid \f{\#e6e5e1}; series colours
\f{\#2a78d6}, \f{\#d6762a}, \f{\#3a9a5b}), \f{period\_num(p)} (year + (quarter $-$ 1)/4) and
\f{save\_fig(p, name, w = 6.3, h)} (PDF to \f{output/study/part1/figures/}). Do not edit
\f{01\_descriptives.R}.

\textbf{Figures} (\f{output/study/part1/figures/}; every figure: axis titles with units, caption-free, sample
named in the document caption).
\par\noindent\begin{tabularx}{\textwidth}{@{}>{\raggedright\arraybackslash}p{0.38\textwidth} Y@{}}
\toprule
File & Content \\
\midrule
\f{national\_totals.pdf} & three panels, free y scales, by quarter: count (billion), amount (Rs.\ lakh crore),
  registered users (million). $h=2.9$. \\
\f{national\_growth.pdf} & bars of \f{qoq\_count} and \f{qoq\_amount} (two panels), 2018Q2--2026Q2; zero line.
  $h=3.0$. \\
\f{avg\_value.pdf} & national \f{atv} line; median and interquartile band of state \f{atv} (main sample). $h=3.2$. \\
\f{type\_shares.pdf} & national \f{share\_*} in percent, three lines with direct labels. $h=3.2$. \\
\f{state\_count\_pc.pdf} & dot plot, 35 rows sorted by 2026Q2 \f{count\_pc}; hollow dot 2018Q4, filled dot
  2026Q2, segment between; log x axis (transactions per person per quarter); y labels are names; the five other
  union territories in grey text. $h=6.5$. \\
\f{reach\_intensity.pdf} & 2026Q2 scatter: x \f{users\_per\_1000}, y \f{intensity} (transactions per user per
  quarter); text labels \f{code}; other union territories in grey. $h=4.0$. \\
\f{decomposition.pdf} & by quarter, main sample: panel 1 \f{var\_y}, \f{var\_r}, \f{var\_i}; panel 2 \f{reach} and
  \f{intensity} contributions. $h=3.2$. \\
\f{retail\_share.pdf} & 2026Q2 scatter: x \f{count\_pc} (log), y \f{share\_retail} (percent); labels \f{code}.
  $h=4.0$. \\
\f{population\_growth.pdf} & bars, 35 ids sorted: percent change of population 2018Q1 to 2026Q2. $h=6.0$. \\
\bottomrule
\end{tabularx}

\textbf{Tables} (\f{output/study/part1/tables/}, through \f{write\_tabular}).
\par\noindent\begin{tabularx}{\textwidth}{@{}>{\raggedright\arraybackslash}p{0.38\textwidth} Y@{}}
\toprule
File & Content \\
\midrule
\f{units.tex} & longtable (caption ``The 35 state units'', label \f{tab:units}): name, \f{state\_id}, code, main
  sample (Yes/No), PhonePe raw names, NCP raw names (from \f{state\_crosswalk.csv} rows that occur in
  \f{data/manual/population\_ncp2020.csv} and in the 2026Q2 PhonePe hover file; joined by ``; ''). \\
\f{checks.tex} & the four rows of \f{phonepe\_checks.csv} with plain-text descriptions, plus two population rows:
  NCP state sum / India for 2018 and for 2027 (7 decimals), and missing population cells in 1190 (0). \\
\f{national\_annual.tex} & \f{annual()}: year, count (billion, 2 dec.), amount (Rs.\ lakh crore, 2), users at
  year end (million, 1), transactions per person (1), average value (Rs., 0), growth of count (percent, 1;
  ``--'' for \f{NA}). \\
\f{state\_2026q2.tex} & 35 rows sorted by \f{count\_pc} descending: name, population (million, 2), transactions
  per person (1), amount per person (Rs., 0), users per 1000 (0), transactions per user (1), average value (Rs., 0),
  retail share (percent, 1). Other union territories marked with a dagger and a note line. \\
\f{concentration.tex} & rows count, amount, users, population; columns: share of the top five (by count) in
  2018Q4 and in 2026Q2 (percent, 1); last row: states needed for half of all transactions. \\
\f{decomposition.tex} & columns 2018Q4, 2026Q2; rows $n$, $\operatorname{var}(y)$, $\operatorname{var}(r)$,
  $\operatorname{var}(i)$, $2\operatorname{cov}(r,i)$, reach contribution, intensity contribution (3 dec.). \\
\bottomrule
\end{tabularx}

\textbf{Numbers} (\f{output/study/part1/numbers.tex}). Families, with \f{<p>} in \{2018Q1, 2018Q4, 2026Q2\}
unless stated. Decimals in brackets. Names are character values.
\begin{longtable}{@{}p{0.42\textwidth} p{0.53\textwidth}@{}}
\toprule
Key & Definition \\
\midrule
\endhead
\f{data.n\_ids}, \f{data.n\_main}, \f{data.n\_other\_ut}, \f{data.n\_variants}, \f{data.n\_quarters},
  \f{data.n\_phonepe\_names} & 35, 30, 5, crosswalk rows, 34, raw hover names in 2026Q2 [0] \\
\f{data.rows.phonepe}, \f{data.rows.types}, \f{data.rows.population} & rows of the interim files [0] \\
\f{check.type\_sum.n}, \f{check.type\_sum.flagged}, \f{check.zero\_fill.n}, \f{check.zero\_fill.flagged},
  \f{check.nat\_count.n}, \f{check.nat\_count.flagged}, \f{check.nat\_users.n}, \f{check.nat\_users.flagged}
  & from \f{phonepe\_checks.csv} [0] \\
\f{pop.ncp\_ratio.2018}, \f{pop.ncp\_ratio.2027} & sum of NCP state rows / India row [7] \\
\f{pop.total.2018Q1}, \f{pop.total.2026Q2}, \f{pop.growth\_pct} & million [1]; percent change [1] \\
\f{pop.growth.max.name}, \f{pop.growth.max}, \f{pop.growth.min.name}, \f{pop.growth.min} & 35 ids, percent
  change 2018Q1--2026Q2 [1] \\
\f{nat.count.<p>}, \f{nat.amount.<p>}, \f{nat.users.<p>} & billion [2], Rs.\ lakh crore [2], million [1] \\
\f{nat.count\_ratio}, \f{nat.amount\_ratio}, \f{nat.users\_ratio} & 2026Q2 / 2018Q1 [0] \\
\f{nat.count\_rate}, \f{nat.amount\_rate}, \f{nat.users\_rate} & \f{annual\_rate}, 2018Q4 to 2026Q2, percent per
  year [1] \\
\f{nat.count\_pc.<p>}, \f{nat.users\_per\_1000.<p>} & national [1], [0] \\
\f{nat.qoq.count.2020Q2}, \f{nat.qoq.amount.2020Q2}, \f{nat.n\_qoq\_neg\_count} & percent [1]; quarters with
  \f{qoq\_count} $<0$ [0] \\
\f{nat.qoq.count.median\_last8} & median \f{qoq\_count} of 2024Q3--2026Q2 [1] \\
\f{nat.n\_states\_fall.2020Q2} & ids (of 35) with a fall in count in 2020Q2 [0] \\
\f{nat.atv.2018Q1}, \f{nat.atv.peak}, \f{nat.atv.peak\_period}, \f{nat.atv.2026Q2} & Rs.\ [0]; period of the
  maximum (character) \\
\f{nat.share.p2p.<p>}, \f{nat.share.retail.<p>}, \f{nat.share.utility.<p>} & percent [1] \\
\f{nat.share.retail.min2018}, \f{nat.share.retail.max2018} & minimum and maximum over 2018Q1--2018Q4 [1] \\
\f{st.count\_pc.max.name.<q>}, \f{st.count\_pc.max.<q>}, \f{st.count\_pc.min.name.<q>}, \f{st.count\_pc.min.<q>},
  \f{st.count\_pc.maxmin.<q>} & main sample, $q \in$ \{2018Q4, 2026Q2\} [1]; max/min ratio [1] \\
\f{st.users\_per\_1000.max.name}, \f{st.users\_per\_1000.max}, \f{st.users\_per\_1000.min.name},
  \f{st.users\_per\_1000.min} & main sample, 2026Q2 [0] \\
\f{st.intensity.max.name}, \f{st.intensity.max}, \f{st.intensity.min.name}, \f{st.intensity.min} & main sample,
  2026Q2 [1] \\
\f{st.rank\_cor} & \f{spearman} of \f{count\_pc} 2018Q4 against 2026Q2, main sample [3] \\
\f{st.reach\_intensity\_cor} & \f{spearman} of \f{users\_per\_1000} and \f{intensity}, 2026Q2, main [3] \\
\f{conc.top5.names.<q>}, \f{conc.top5.count.<q>}, \f{conc.top5.population.<q>}, \f{conc.n\_half.<q>} & names
  joined by ``, '' and ``and''; percent [1]; integer \\
\f{dec.<q>.var\_y}, \f{dec.<q>.var\_r}, \f{dec.<q>.var\_i}, \f{dec.<q>.reach}, \f{dec.<q>.intensity} & [3] \\
\f{dec.var\_y.max}, \f{dec.var\_y.max\_period} & maximum of \f{var\_y} over quarters [3]; its period \\
\f{retail.cor.2026Q2} & \f{spearman} of \f{share\_retail} and \f{count\_pc}, 2026Q2, main [3] \\
\f{retail.max.name}, \f{retail.max}, \f{retail.min.name}, \f{retail.min} & main sample, 2026Q2, percent [1] \\
\bottomrule
\end{longtable}
Each entry has a \f{source} comment: the function and the input (for example
\f{national(); data/interim/phonepe.csv}). The prose tasks may need more keys; they add them in this script
with the same rules and a test-free review.

\begin{steps}
\item Write \f{study\_theme.R} and \f{part1\_study.R}. The script sources the three libraries, reads the
      inputs with \f{load\_part1}, builds the figures, tables and numbers above, and calls
      \f{write\_numbers(..., generator = "Rscript src/R/part1\_study.R")}.
\item Run \f{Rscript src/R/part1\_study.R}. Expected: exit 0; 9 PDF files; 6 \f{.tex} tables;
      \f{grep -c defval output/study/part1/numbers.tex} is at least 100.
\item Check values against the scratch run in the spec: \f{dec.2018Q4.reach} is 0.776, \f{dec.2026Q2.intensity}
      is 0.614, \f{st.rank\_cor} is 0.892, \f{nat.count.2026Q2} is 38.66, \f{check.zero\_fill.flagged} is 1.
      A difference means a defect in the code; find it before you continue.
\item Look at each PDF (Read tool). Check: no clipped labels, no overlapping text that hides a value, axis units
      present. Fix and run again.
\item Run the script a second time and \f{git diff --exit-code output/study/part1/tables output/study/part1/numbers.tex}
      after a commit of the first run (determinism). Expected: exit 0.
\item Run \f{Rscript -e 'testthat::test\_dir("tests/R")'}. Expected: all PASS.
\item Commit the two R files, \f{Makefile} and \f{output/study/part1/}:
      \f{git commit -m "feat: Part 1 study figures, tables and computed numbers"}.
\end{steps}

% =====================================================================
\section{Task 6: Study document skeleton}
\textbf{Files.} Create \f{thesis/study/part1.tex}. Modify \f{Makefile} (\f{report}: add
\verb|cd thesis/study && for f in part*.tex; do latexmk -pdf -interaction=nonstopmode $$f || exit 1; done|).
Modify \f{.gitignore} only if needed so that \f{thesis/study/part1.pdf} is tracked (it is tracked today: only
\f{thesis/btp1/main.pdf} is ignored).

\textbf{Preamble.} As \f{thesis/progress/part3.tex}, plus \f{tikz} (library \f{arrows.meta, positioning}),
\f{longtable}, \f{caption} (\f{font=small}). \verb|\graphicspath{{../../output/study/part1/figures/}}|. After
\f{hyperref}:
\begin{verbatim}
\makeatletter
\newcommand{\defval}[2]{\expandafter\def\csname val@#1\endcsname{#2}}
\newcommand{\val}[1]{\ifcsname val@#1\endcsname\csname val@#1\endcsname\else
  \PackageError{study}{Unknown value key `#1'}{Run Rscript src/R/part1_study.R}\fi}
\makeatother
\input{../../output/study/part1/numbers.tex}
\end{verbatim}
Title: ``Digital payment adoption across Indian states, 2018--2026: building the state panel from PhonePe Pulse''
with the line ``BTP study, Part 1''; author Devansh Gupta; date November 2026.

\textbf{Sections} (labels in brackets): Abstract; 1 Introduction [\f{sec:intro}]; 2 Data sources [\f{sec:data}];
3 Construction of the panel [\f{sec:construction}]; 4 Adoption in India as a whole [\f{sec:national}];
5 Adoption across states [\f{sec:states}]; 6 Limitations of the data [\f{sec:limits}]; 7 Conclusions
[\f{sec:conclusions}]; References; Appendix A The 35 state units [\f{app:units}]; Appendix B Reproduction
[\f{app:repro}]. Subsections as in the spec Section 4.4.

\begin{steps}
\item Create the file with the preamble, all section and subsection headings, and every float: Figure
      \f{fig:flow} (TikZ: boxes for the three raw sources $\to$ the three Python modules $\to$ the four interim
      files $\to$ \f{part1\_study.R} $\to$ figures, tables, numbers); figures \f{fig:<file>} for the nine PDFs;
      tables \f{tab:files} (data files used: hand-written \f{tabular}, file counts by \verb|\val|),
      \f{tab:checks}, \f{tab:annual}, \f{tab:states}, \f{tab:concentration}, \f{tab:decomposition}
      (\verb|\input| of the generated files), and \verb|\input| of \f{units.tex} in Appendix A. Each float has its
      final caption (unit, sample, period, source with \verb|\citep|). Under each heading, one \verb|%| comment
      line says what the section will contain (from spec Section 4.4).
\item Write Appendix B now: the command sequence \f{make data analysis report}, the pinned commit, the input
      files, and the statement that \f{numbers.tex} is generated and not edited.
\item Run \f{cd thesis/study \&\& latexmk -pdf -interaction=nonstopmode part1.tex}. Expected: exit 0;
      \f{grep -E "undefined|Unknown value key" part1.log} prints nothing; no \f{Overfull} line wider than 1\,pt.
\item Negative check: add \verb|\val{no.such.key}| to the body, compile, see the error \f{Unknown value key},
      remove the line.
\item Commit \f{thesis/study/part1.tex}, \f{part1.pdf}, \f{Makefile}:
      \f{git commit -m "docs: Part 1 study skeleton with figures and tables"}.
\end{steps}

% =====================================================================
\section{Task 7: Prose, Sections 1--3}
\textbf{Files.} Modify \f{thesis/study/part1.tex}; if a new key is needed, \f{src/R/part1\_study.R} and the
regenerated \f{output/study/part1/numbers.tex}.

\textbf{Content.}
\begin{description}
\item[1 Introduction] (4--6 paragraphs). The BTP question (adoption and financial inclusion across states,
  \f{proposal.txt}). Why a state panel over time is necessary. NPCI state data only from April 2025
  (cite \f{npci\_upi\_state}, \f{medianama2025npci}); therefore PhonePe Pulse (\f{phonepe\_pulse}), used before by
  Reddy, Kedia and Shukla (\f{reddy2024icrier}, with \verb|\citet|). The aim of this part: build the state-quarter panel of PhonePe adoption and its
  population denominator, check it, and describe it in nominal terms. The contents of each section. One sentence
  each on Part 2 (price deflation, comparison with NPCI, maps) and Part 3 (convergence analysis).
\item[2 Data sources] PhonePe Pulse: what it is; licence CDLA-Permissive-2.0; the pinned commit and its date;
  the three folders; the restatement of the series from 2018Q1 and the README rule against joins with earlier
  releases; the files used and their content (Table \f{tab:files}); the type definitions; the files not used and
  why (spec Section 4.2). NCP 2020 (\f{ncp2020}): the report, Table 8, 1 March reference date, thousands,
  years 2018--2027, cohort-component projections from the 2011 Census. Figure \f{fig:flow}.
\item[3 Construction] 3.1 State units: the crosswalk, its name variants (key \f{data.n\_variants}), the 35 ids, merged ids
  (Ladakh into Jammu and Kashmir; Dadra and Nagar Haveli and Daman and Diu), the main sample and the five other
  union territories; reference to Appendix A. 3.2 Parsing and checks: the summation rule for merged ids; the
  stop rules; Table \f{tab:checks} with the numbers; the zero-filled entry. 3.3 Population: the formula
  $P_{y,q} = P_y + \frac{q-1}{4}(P_{y+1}-P_y)$, why 2027 is needed, the NCP sum ratios, Figure
  \f{fig:population\_growth} with the highest and lowest growth. 3.4 Measures: the formulas of spec Section 4.3,
  each with its unit and meaning, and the decomposition identity.
\end{description}

\begin{steps}
\item Write the prose. Every number is \verb|\val{...}| or in a table. Each figure and table is introduced before
      it appears and explained after (what it shows, then what the numbers mean).
\item Compile as in Task 6. Expected: same checks pass.
\item Run the banned-word check:
\begin{verbatim}
grep -nwiE "very|significant(ly)?|robust|dramatic(ally)?|huge|clearly|obviously|interestingly|
  notably|crucial|leverage|boost|massive|remarkabl[ey]|impressive|explosive|journey|landscape|
  delve|student|AI|Task" thesis/study/part1.tex | grep -v '^[0-9]*:\s*%'
\end{verbatim}
      (one line in the shell). Expected: no output, except words inside citation keys.
\item Read the PDF pages of Sections 1--3. Check: sentences of 25 words or fewer in descriptive text; active voice;
      no claim of a cause without evidence.
\item Commit \f{part1.tex}, \f{part1.pdf} (and the script and \f{numbers.tex} if changed):
      \f{git commit -m "docs: Part 1 study, introduction, data and construction"}.
\end{steps}

% =====================================================================
\section{Task 8: Prose, Sections 4--7 and the abstract}
\textbf{Files.} As Task 7.

\textbf{Content.}
\begin{description}
\item[4 India as a whole] Figure \f{fig:national\_totals} and Table \f{tab:annual}: levels and ratios
  2018Q1--2026Q2, annual rates 2018Q4--2026Q2. Figure \f{fig:national\_growth}: the 2020Q2 fall (the only quarter
  with a fall in the national count; number of states with a fall), the lower growth in the last eight quarters.
  Figure \f{fig:avg\_value}: rise to the peak, then fall; state band. Figure \f{fig:type\_shares}: the 2018 swing,
  then the rise of the retail share. Interpretation only where the data show it: for example, a lower average
  value at the same time as a higher retail share is consistent with more small merchant payments; the data do
  not give the value by type, so the text says that this is not tested.
\item[5 Across states] Figure \f{fig:state\_count\_pc} and Table \f{tab:states}: highest and lowest, max/min
  ratios in 2018Q4 and 2026Q2, rank persistence. Table \f{tab:concentration}: top five against their population
  share; states for half of all transactions. Figure \f{fig:reach\_intensity}: reach and intensity, extremes,
  their rank correlation. Figure \f{fig:decomposition} and Table \f{tab:decomposition}: reach explained most of
  the cross-state variance in 2018Q4, intensity in 2026Q2 (use the values); the peak of \f{var\_y}. Figure
  \f{fig:retail\_share}: the rank correlation and the extremes.
\item[6 Limitations] One paragraph each: one app, with its share of UPI a Part 2 question; registered users are a
  cumulative count (no fall in any state-quarter is consistent with this) and not active users; state attribution
  method not documented in the repository; projections, not counts; nominal values; type shares in 2018.
\item[7 Conclusions] Four to six findings, each with its number, and what Parts 2 and 3 take from the panel.
\item[Abstract] 150--250 words: data, method, the main numbers.
\end{description}

\begin{steps}
\item Write the prose with the rules of Task 7.
\item Compile, banned-word check and PDF read as in Task 7.
\item Check the page count: \f{pdfinfo thesis/study/part1.pdf | grep Pages}. Expected: 20 to 30.
\item Commit: \f{git commit -m "docs: Part 1 study, results, limitations and conclusions"}.
\end{steps}

% =====================================================================
\section{Task 9: Full rebuild, checks and tag}
\textbf{Files.} No new files expected. Fixes go to the file that owns the defect.

\begin{steps}
\item Run \f{make test}. Expected: pytest all pass; R tests 0 failures.
\item Run \f{make data analysis report}. Expected: exit 0.
\item Run \f{git status --short}. Expected: only figure PDFs (creation date) and LaTeX PDFs show as changed. Run
      \f{git diff --exit-code data/ output/tables output/study/part1/tables output/study/part1/numbers.tex}.
      Expected: exit 0. Restore the changed PDFs that are not part of this work with \f{git checkout --}
      (\f{output/figures/*.pdf}, \f{thesis/btp1-slides/slides.pdf}, \f{thesis/progress/*.pdf}).
\item Check every \verb|\val| key in the document exists:
\begin{verbatim}
comm -23 <(grep -o '\\val{[^}]*}' thesis/study/part1.tex | sed 's/\\val{\(.*\)}/\1/' | sort -u) \
         <(grep -o '^\\defval{[^}]*}' output/study/part1/numbers.tex | sed 's/\\defval{\(.*\)}/\1/' | sort -u)
\end{verbatim}
      Expected: no output (the compile also enforces this).
\item Check the document has no reference to ``student'', ``AI'', ``Task'' or ``progress report''
      (\f{grep -niwE "student|AI|Task|progress report" thesis/study/part1.tex}). Expected: no match outside comments.
\item Commit the rebuilt \f{thesis/study/part1.pdf} if it changed:
      \f{git commit -m "build: rebuild Part 1 study"}.
\item Create the tag (do not push): \f{git tag -a part1-study -m "Part 1 study: PhonePe Pulse state panel, 2018--2026"}.
      Check: \f{git tag -n | grep part1}. Expected: \f{part1} (unchanged) and \f{part1-study}.
\end{steps}

\end{document}
